% Independent review fragment; amsmath, amssymb and hyperref are sufficient.
\section{Independent review of the QN5 and QN6 logical composition}

The quantifier-head exchange, formula radix change, and no-push MP controller
are compatible. This review is an ordinary mathematical argument using the
supplied syntactic PA calibration. It is not a new end-to-end Lean or Rocq
theorem. Candidate-specific source bindings, positive primitive simulations,
and complete blank startup are separate obligations.

\subsection{Precisely which controller is being compiled}
Let $p(a,b)=(a+b)(a+b+1)/2+b$, $[a;w]=p(a,w)$,
$I_k(a,b)=kp(a,b)$, and $U_k(i,a)=kp(a,i)+2$, with $k=5$ or $6$.
Stores are $(t,w,d,h)$. QN retains the original Shift guard:
\[
 S(t,w,d,h)=\begin{cases}\mathsf{halt}&t+h=1,\\
 (h,\operatorname{tl}w,d,\operatorname{hd}w)&t+h\ne1.
 \end{cases}
\]
With $v=d+x$, its MP rule is
\[
 \operatorname{MP}_x(t,w,d,h)=
 \begin{cases}(v,w,v,0)&I_k(t,v)=h,\\(0,w,v,0)&I_k(t,v)\ne h.
 \end{cases}
\]
It does not push a dummy zero on success. It does not return the conclusion
in the cached slot. Both distinctions matter for the MP proof word below.

The AX output is
\[
 (g,[q_1,\ldots,q_9,t;w],0,0),
\]
where $q_1$ is the diagonal axiom, $q_2,\ldots,q_7$ are the six logical
heads in reverse order, $q_8$ is the dominated implication head, and
$q_9=U_k(d,h)$. The Q exchange changes only $q_2$.
The gate and $q_1,\ldots,q_8$ are multiples of $k$, whereas
$q_9\equiv2\pmod k$ and $q_9\ge2$. These constructor identities hold
for every numerical parameter, independently of consistency. Consequently
\[
 \operatorname{AX};S^9=(q_8,[t;w],0,U_k(d,h)),\qquad
 \operatorname{AX};S^{10}=(U_k(d,h),w,0,t),
\]
with the appropriate unchanged AX parameters. No displayed Shift halts.
The argument uses the stated residue classes and positivity, not a special
identity of radix four.

\subsection{Exact stack words and the reserved zero}
Put $A_0=\operatorname{AX}_{0,0,0}$ and compose words left to right. Then
\[
 \operatorname{Push}=A_0;S^9;\operatorname{MP}_0,
 \qquad \operatorname{Pop}=A_0;S^{11},
 \qquad \operatorname{Discard}=\operatorname{MP}_0,
 \qquad \operatorname{Offset}_i=\operatorname{MP}_i.
\]
Their contracts are
\begin{align*}
 \operatorname{Push}(t,w,d,h)&=(0,[t;w],0,0),\\
 \operatorname{Pop}(t,[a;w],d,h)&=(t,w,0,a),\\
 \operatorname{Discard}(t,w,0,h)&=(0,w,0,0),\\
 \operatorname{Offset}_i(0,w,0,0)&=(0,w,i,0).
\end{align*}
Push's MP comparison fails because implication and universal codes have
different residues. Pop's eleventh Shift tests $U_k(d,h)+t\ge2$, so Pop
is total even when $t=1$. Discard's two MP branches coincide when $d=x=0$,
even with an arbitrary old cache. Offset uses $I_k(0,i)>0$ for $i>0$ and
the coincident zero outcomes for $i=0$.

For $1\le j\le7$, the head selector
\[
 \operatorname{Discard};\operatorname{AX}_{A,D,H};
 S^{j+1};\operatorname{Pop}^{9-j}
\]
returns $(q_j,w,0,0)$ from $(t,w,0,0)$. The last Pop consumes the saved
zero. For the dominated head with $u\le D$, use
\[
 \operatorname{Discard};\operatorname{Offset}_u;
 \operatorname{AX}_{A,D-u,H};S^9;\operatorname{Pop}.
\]
The gate selector is
$\operatorname{Discard};\operatorname{AX}_{c,b,a};\operatorname{Pop}^{10}$.
Generalization uses
\[
 \operatorname{Push};\operatorname{Offset}_i;S;
 \operatorname{AX}_{i,0,0};S^{10}
\]
and returns $(U_k(i,a),w,0,0)$ from $(a,w,0,0)$.
These words remain valid for Q because the only changed head is selected
at the same position and has the same residue.

For proofs of $a$ and $I_k(a,b)$, the correct no-push MP compilation is
\[
 \operatorname{Discard};\operatorname{Push};
 C(I_k(a,b));\operatorname{Push};C(a);\operatorname{Push};
 S;S;\operatorname{MP}_b;\operatorname{Pop}.
\]
If neither recursive call halts, the three pushes yield
$(0,[a,I_k(a,b),0;w],0,0)$. The first Shift yields
$(0,[I_k(a,b),0;w],0,a)$. The second halts if $a=1$; otherwise it yields
$(a,[0;w],0,I_k(a,b))$. MP succeeds, giving $(b,[0;w],b,0)$, and Pop
returns $(b,w,0,0)$. This is why the compiler must reserve a zero beneath
the premise computations. The old controller's suffix that expects MP to
push its own zero must not be reused.

Every identity uses the actual Cantor projections, including
$[0;0]=0$ and $\operatorname{hd}(0)=\operatorname{tl}(0)=0$.
No free-list representation or nonempty-stack hypothesis is introduced.
Induction on a proof gives a primitive-recursive word $C(d)$ which either
halts or returns its conclusion on the original arbitrary stack, with
$d=h=0$. For a refutation, $A_0;C(d);S$ halts from every incoming store.
This uniformity is required because the scheduler carries the store from
one trial to the next.

\subsection{The arithmetic calibration and its delicate points}
The inherited radix translation must recursively recode both formulas and
variable labels. For $n>0$, writing $n=4p(a,b)+r$ with $r<4$, set
\[
 F_k(0)=0,\qquad F_k(n)=kp(F_k(a),F_k(b))+r.
\]
The smaller-argument recursion is primitive recursive. Its stated left
inverse proves injectivity, preserving label equalities, disequalities,
binding and shadowing. In particular,
\[
 F_k(I_4(B,U_4(B,B)))=I_k(F_k(B),U_k(F_k(B),F_k(B))).
\]
Keeping labels fixed while recoding formulas would invalidate this diagonal
identity and is not part of the proof.

Numerical order is not preserved: $40<44$ but
$F_5(40)=75>55=F_5(44)$. Restricted propositional leaves must therefore be
expanded using the finite generic \texttt{rho8All} proof, instantiated
with numerical $\le$ and $I_k$. The required closure bounds follow from
$a,b\le p(a,b)\le I_k(a,b)$. The generic rules in the supplied
\texttt{RestrictedImplication.lean} require only MP, the restricted head,
reflexivity, and these two operand bounds; they do not require injective
constructors, the exchanged quantifier head, or unrestricted substitution
through a guarded gate.

The same generic basis proves exchange
$(P\to(Q\to R))\to(Q\to(P\to R))$ at radix $k$.
Applying it in both orientations gives conclusion-preserving
primitive-recursive proof translations between the old sixth head and
$U_k(D,I_k(A,H))\to(U_k(D,A)\to U_k(D,H))$.
This supplies Q's logical equivalence independently of the controller.

For the reverse arithmetic translation, residues $4,\ldots,k-1$ are
interpreted as closed truth. They cannot be omitted, since arbitrary
command parameters can contain them. The alias $I_k(0,0)=0$ is handled
by the fixed proof that truth is equivalent to truth implying truth.
Code one is interpreted as the irreflexive bit-membership assertion and
has a uniform PA refutation. In the guarded gate, the supplied finite
comprehension proof must distinguish $a=b$ from $a\ne b$, preserve
capture-avoiding substitution, and use freshness of label $c$ in the
translation of formula code $c$. The supplied radix argument makes all
three distinctions and uses a syntactically generated PA induction
instance, not a satisfaction predicate or reflection principle.

Accordingly the proof chain is
\[
 \mathrm{PA\ refutation}\longrightarrow\mathcal C_4\mathrm{\ refutation}
 \longrightarrow\mathcal C_k\mathrm{\ refutation}
 \longrightarrow\mathcal C_k^{\mathrm Q}\mathrm{\ refutation},
\]
with the reverse direction obtained by exchange followed by the total
radix-$k$ arithmetic translation. One compiles a fresh target-radix
command word after translating a proof. Neither addition of offsets nor
ordinary Cantor stack pairing is conjugated by $F_k$.

\subsection{Source and proof boundary}
Replacing the four-increment Cons body by $k$ increments changes its loop
to $(T,Y)\mapsto(T+k,Y-1)$ on the positive branch. Its invariant is
$T=T_0+ki$, $Y=Y_0-i$, and its exact cost is $(k+2)Y_0+1$ primitive
steps. The no-push splice removes the earlier 19-site zero-push loop;
the unchanged retained copy loop transfers the consequent to $T$ and $D$.
These modifications occur in different source subtrees. The candidate
binding checks must retain their exact CFG changes, all 61 Cons occurrences,
and all external successor edges; a four-increment lemma alone would not
prove the changed arithmetic routine.

The sound-store invariant annotates $T,H$ and every canonical stack entry
with source proofs; offset $D$ has no proof obligation. AX preserves the
annotations through its axiom suppliers and generalization; successful MP
combines its annotated premises, failure returns truth, and Shift transports
the annotations. At a Shift halt, $t+h=1$ exposes a proof of code one.
The all-store refutation word and the cofinal scheduler supply the converse.
After the independently checked positive physical simulation and complete
blank startup, these ordinary PA-formalizable arguments yield the requested
PA-provable equivalence for the exact QN5 and QN6 tables. They do not close
the repository's separate iPRA or end-to-end proof-assistant endpoint.
